\chapter{Python}
\label{chap:python}
\vspace{10.5cm}
%~ foo
%~ \newpage{}
\section{PEP}
\section{Built-in Data TYpes}
\subsection{\href{https://docs.python.org/3/library/stdtypes.html\#dict}{dict}}
\subsection{\href{https://docs.python.org/3/library/stdtypes.html\#list}{list}}
\subsection{\href{https://docs.python.org/3/library/stdtypes.html\#set}{set}}
\subsection{\href{https://docs.python.org/3/library/stdtypes.html\#frozenset}{frozenset}}
\subsection{\href{https://docs.python.org/3/library/stdtypes.html\#tuple}{tuple}}
\subsection{\href{https://docs.python.org/3/library/stdtypes.html\#string}{string}}
\subsection{\href{https://docs.python.org/3/library/stdtypes.html\#bytes}{bytes}}
\subsection{\href{https://docs.python.org/3/library/stdtypes.html\#bytearray}{bytearray}}

\newpage
\section{Decorator}
A function decorator is a callable object that takes a function as its 
input parameter
\href{https://www.programiz.com/python-programming/decorator}{decorator} \\
\href{https://www.geeksforgeeks.org/decorators-with-parameters-in-python/}{Decorators with parameters} \\
\href{https://wiki.python.org/moin/PythonDecoratorLibrary}{PythonDecoratorLibrary}
\newpage{}
\section{Special Instance Method}
\href{https://docs.python.org/3/reference/datamodel.html#specialnames}{Special Instance Method}


\section{Providing Multiple Ctors}
\href{https://realpython.com/python-multiple-constructors/}{multiple \texttt{\char`_\char`_init\char`_\char`_}}

\newpage
